Tento výzkumný úkol se zabývá teorií a návrhem propojení umělé inteligence (AI) a strojového učení ve hrách k dynamickému přizpůsobování chování protivníků v reálném čase.

V první kapitole se popisuje umělá inteligence ve hrách. Specificky popisuje historii umělé inteligence ve stolních a digitálních hrách, dále rozebírá některé algoritmy, které jsou nejčastěji používané k vytvoření umělé inteligence v digitálních hrách, ale díky nim není umělá inteligence schopna se přizpůsobit nebo upravit své chování. Nakonec se tato část specificky zaměří na některé existující 3D akční hry, které obsahují inteligentnější AI.

Ve druhé kapitole se zmiňují základy strojového učení a také existující čtyři typy strojového učení:
\begin{enumerate}
	\item \uv{Učení s učitelem}
	\item \uv{Učení bez učitele}
	\item \uv{Částečně řízené učení}	
	\item \uv{Zpětnovazebné učení}
\end{enumerate}
A také jeden nebo dva algoritmy z každého typu strojového učení.

Kapitola třetí obsahuje návrh autonomního agenta. Především obsahuje návrh samotného prostředí, kde se agent bude učit a zároveň testovat svou inteligenci. Dále tato kapitola obsahuje dva algoritmy, SARSA a Deep Q-Learning, které budou použity pro sestrojení agenta pro 3D akční hru.

Ve čtvrté kapitole řeším typy hodnocení, podle kterých by se agent rozhodoval, v čem se zlepšit a co nepotřebuje využívat. Jsou tu tři návrhy hodnocení, které by představovaly zajímavou obtížnost pro agenta, aby se jim přizpůsobil. 

Závěr obsahuje konečné myšlenky autora a také pokračování pro další práci, kde by se myšlenka tohoto výzkumného úkolu mohla uskutečnit.